\documentclass[12pt]{article}
\usepackage[utf8]{inputenc}
\usepackage[T1]{fontenc}
\usepackage[margin=1in]{geometry}
\usepackage{setspace}
\usepackage{xcolor}
\usepackage[colorlinks=true, urlcolor=blue!60!black, linkcolor=blue!60!black]{hyperref}
\setlength{\parindent}{0em}
\setlength{\parskip}{0.6em}
\renewcommand{\baselinestretch}{1.0}
\begin{document}
\begin{center}
{\LARGE \textit{Wickard v. Filburn}, 317 U.S. 111 (1942)\par}
\vspace{0.5em}
{\large Case Brief --- Two-Page Edition\par}
\end{center}

\section{Memory Jogger}
Congress may regulate home-grown, home-consumed wheat because, aggregated across many farmers, it substantially affects interstate commerce (\href{https://supreme.justia.com/cases/federal/us/317/111/}{317 U.S. 111}).

\section{Detailed Case Facts}
The Agricultural Adjustment Act of 1938 set farm-by-farm wheat allotments and marketing quotas to stabilize prices. Ohio farmer Roscoe Filburn, allotted 11.1 acres, sowed 23 and used much of his wheat for feed, flour, and seed. A 1941 amendment raised the penalty on excess wheat to 49 cents a bushel and extended it to wheat consumed on the farm; Filburn refused to pay his \$117.11 penalty and sued Secretary Wickard.

\section{Procedural History}
A divided three-judge district court enjoined enforcement; after reargument on the commerce question, the Supreme Court reversed.

\section{Judicial Votes}
9--0. Jackson wrote for a unanimous Court; no separate opinions.

\section{Holding}
Local activity that is not commerce may be regulated if it exerts a substantial economic effect on interstate commerce, whether direct or indirect, and trivial individual contributions are aggregated with those of others similarly situated. The penalty did not violate due process.

\section{Analysis of Opinions}
Jackson discarded the formal production-versus-commerce categories of \textit{E.\ C.\ Knight} and \textit{Carter Coal} for a practical economic-effects test. Home-grown wheat, the most variable factor in wheat consumption, overhangs the market and displaces purchases, so it competes with wheat in commerce. Since cooperating farmers received subsidies, regulating them was not arbitrary, and the penalty attached only on threshing after the amendment.

\section{Examples: Future Applications}
\textbf{Same-Side:} home-grown medical marijuana may be federally banned (\textit{Raich}). \textbf{Opposite-Side:} gun possession near schools cannot be aggregated (\textit{Lopez}). \textbf{Fence-Sitter:} a federal ban on home distilling for personal use.

\section{Critique}
Barnett and Epstein argue that \textit{Wickard} abandoned the original meaning of commerce and, with aggregation, erased enumerated limits; defenders see a necessary response to national collective-action problems that later cases have bounded.

\section{Key Quotations}
\begin{quote}
``But even if appellee's activity be local, and though it may not be regarded as commerce, it may still, whatever its nature, be reached by Congress if it exerts a substantial economic effect on interstate commerce, and this irrespective of whether such effect is what might at some earlier time have been defined as `direct' or `indirect.' ''
\end{quote}
\noindent --- Justice Jackson, Opinion of the Court, 317 U.S. at 125.

\end{document}
